%============================================================
% 03_summaries.tex — 1·2페이지 표 내용
%
% 이 파일의 각 \newcommand가 다음 칸을 채웁니다.
%   1페이지(연구결과 요약문): summaryGoal / summaryResult / summaryEffect
%   2페이지(과제 수행 과정): summaryExpert / summaryProblemSolving /
%                            summaryFirstMember / summarySecondMember / summaryLearning
%
% \jiwon[숫자]는 한글 더미 텍스트입니다 (Lorem ipsum의 한글 버전).
% 각 칸의 \jiwon[...] 줄을 지우고 실제 내용을 작성하세요.
% 여러 문단을 써도 되고, 줄바꿈도 사용할 수 있습니다.
%============================================================

%%% 1페이지 ─────────────────────────────────────────────────

\newcommand{\summaryGoal}{%
  % '연구 필요성 및 목적' 칸
  \jiwon[1] % ← 이 줄을 지우고 내용을 작성하세요
}

\newcommand{\summaryResult}{%
  % '연구내용 및 연구결과' 칸
  \jiwon[2-3] % ← 이 줄을 지우고 내용을 작성하세요
}

\newcommand{\summaryEffect}{%
  % '연구 기대효과' 칸
  \jiwon[4] % ← 이 줄을 지우고 내용을 작성하세요
}

%%% 2페이지 ─────────────────────────────────────────────────

\newcommand{\summaryExpert}{%
  % '전문가 피드백 활용 과정' 칸
  % 전문가 피드백을 받지 않았다면 아래 한 줄로 대체하세요.
  %   해 당 없 음
  \jiwon[5] % ← 이 줄을 지우고 내용을 작성하세요
}

\newcommand{\summaryProblemSolving}{%
  % '문제해결 과정' 칸 — 팀 전체의 과정을 서술
  \jiwon[6-7] % ← 이 줄을 지우고 내용을 작성하세요
}

\newcommand{\summaryFirstMember}{%
  % '팀원1' 칸 — 00_info.tex의 \myauthorfirstko(저자1)의 역할
  \jiwon[8] % ← 이 줄을 지우고 내용을 작성하세요
}

\newcommand{\summarySecondMember}{%
  % '팀원2' 칸 — 00_info.tex의 \myauthorsecondko(저자2)의 역할
  \jiwon[8] % ← 이 줄을 지우고 내용을 작성하세요
}

\newcommand{\summaryLearning}{%
  % '과제 수행 과정에서의 학습' 칸
  \jiwon[9] % ← 이 줄을 지우고 내용을 작성하세요
}
